\documentclass{article}
\usepackage{mathtools} 
\usepackage{fontspec}
\usepackage[UTF8]{ctex}
\usepackage{amsthm}
\usepackage{mdframed}
\usepackage{xcolor}
\usepackage{amssymb}
\usepackage{amsmath}


% 定义新的带灰色背景的说明环境 zremark
\newmdtheoremenv[
  backgroundcolor=gray!10,
  % 边框与背景一致，边框线会消失
  linecolor=gray!10
]{zremark}{说明}

% 通用矩阵命令: \flexmatrix{矩阵名}{元素符号}{行数}{列数}
\newcommand{\flexmatrix}[4]{
  \[
  #1 = \begin{pmatrix}
    #2_{11}     & #2_{12}     & \cdots & #2_{1#4}   \\
    #2_{21}     & #2_{22}     & \cdots & #2_{2#4}   \\
    \vdots      & \vdots      & \ddots & \vdots     \\
    #2_{#31}    & #2_{#32}    & \cdots & #2_{#3#4}
  \end{pmatrix}
  \]
}

% 简化版命令（默认矩阵名为A，元素符号为a）: \quickmatrix{行数}{列数}
\newcommand{\quickmatrix}[2]{\flexmatrix{A}{a}{#1}{#2}}

\begin{document}
\title{习题 1}
\author{张志聪}
\maketitle

\section*{10}

题目是让证明：
\begin{align*}
  \|A\|_P = \max\limits_{x \neq 0} \frac{\|Ax\|_P}{\|x\|_P} = \|P A P^{-1}\|
\end{align*}
因为
\begin{align*}
  \|A\|_P 
  & = \max\limits_{x \neq 0} \frac{\|Ax\|_P}{\|x\|_P} \\
  & = \max\limits_{x \neq 0} \frac{\|PAx\|}{\|Px\|} 
\end{align*}
令
\begin{align*}
  y = Px
\end{align*}
于是可得
\begin{align*}
  x = P^{-1} y
\end{align*}
代入后
\begin{align*}
  \max\limits_{x \neq 0} \frac{\|PAx\|}{\|Px\|} 
  & = \max\limits_{y \neq 0} \frac{\|PAP^{-1}y\|}{\|y\|} \\
  & = \|PA P^{-1}\|
\end{align*}
综上
\begin{align*}
  \|A\|_P = \|P A P^{-1}\|
\end{align*}

\end{document}